\section{Dominated scalar Hermitian normal form}
\label{sec:setup}

Let \(V\) be a finite index set.  Let \(P=(p_{ij})_{i,j\in V}\) be
real symmetric, entrywise nonnegative, and zero on the diagonal.  Let
\(H=H^*\in\C^{V\times V}\) be zero on the diagonal and satisfy
\begin{equation}\label{eq:domination}
 |H_{ij}|\le p_{ij}.
\end{equation}
Write
\[
 q_{ij}=|H_{ij}|,
 \qquad
 H_{ij}=q_{ij}\phi_{ij}
\]
on every edge with \(q_{ij}>0\), where
\[
 |\phi_{ij}|=1,
 \qquad
 \phi_{ji}=\overline{\phi_{ij}}.
\]
The value of \(\phi_{ij}\) on a zero \(q\)-edge is irrelevant.

\begin{definition}[Row support graph]\label{def:row-graph}
The undirected graph \(G_P\) has vertex set \(V\) and edge
\(\{i,j\}\) exactly when \(p_{ij}>0\).  The graph \(G_q\) is defined
similarly using \(q_{ij}>0\).
\end{definition}
These are exact post-aggregation graphs.  In particular, two atomic
contributions that cancel in \(H_{ij}\) remove the edge from \(G_q\),
even if they were separately nonzero before summation.

Assume first that \(G_P\) is connected and \(P\ne0\).  Perron--Frobenius
theory \cite{BermanPlemmons1994,HornJohnson2013} gives
\begin{equation}\label{eq:perron}
 Pr=\rho r,
 \qquad
 r_i>0,
 \qquad
 \sum_i r_i^2=1,
\end{equation}
where \(\rho=\rho(P)\).

\begin{definition}[Two spectral sectors]\label{def:sectors}
For \(\sigma\in\{+1,-1\}\), define
\begin{align}
 \delta_\sigma(P,H)
 &=
 \rho(P)-\lambda_{\max}(\sigma H).
 \label{eq:sector-defect}
\end{align}
Thus
\[
 \delta_+=\rho-\lambda_{\max}(H),
 \qquad
 \delta_-=\rho+\lambda_{\min}(H),
\]
and
\begin{equation}\label{eq:norm-defect}
 \rho(P)-\|H\|=\min\{\delta_+,\delta_-\}.
\end{equation}
\end{definition}

The sign \(\sigma\) is global.  It is not an independently chosen sign
on each edge.

\begin{definition}[Cycle holonomy]\label{def:holonomy}
For an oriented cycle
\[
 C=(v_0,v_1,\ldots,v_\ell=v_0)
\]
contained in \(G_q\), set
\begin{equation}\label{eq:holonomy}
 \Hol_\phi(C)
 =
 \prod_{k=0}^{\ell-1}\phi_{v_kv_{k+1}},
 \qquad
 \Hol_\sigma(C)
 =
 \sigma^\ell\Hol_\phi(C).
\end{equation}
Reversing orientation conjugates the holonomy, so the condition that it
equal \(1\) is orientation independent.
\end{definition}
